% Options for packages loaded elsewhere
\PassOptionsToPackage{unicode}{hyperref}
\PassOptionsToPackage{hyphens}{url}
\documentclass[
  11pt,
  letterpaper,
]{article}
\usepackage{xcolor}
\usepackage[margin=1in]{geometry}
\usepackage{amsmath,amssymb}
\setcounter{secnumdepth}{-\maxdimen} % remove section numbering
\usepackage{iftex}
\ifPDFTeX
  \usepackage[T1]{fontenc}
  \usepackage[utf8]{inputenc}
  \usepackage{textcomp} % provide euro and other symbols
\else % if luatex or xetex
  \usepackage{unicode-math} % this also loads fontspec
  \defaultfontfeatures{Scale=MatchLowercase}
  \defaultfontfeatures[\rmfamily]{Ligatures=TeX,Scale=1}
\fi
\usepackage{lmodern}
\ifPDFTeX\else
  % xetex/luatex font selection
\fi
% Use upquote if available, for straight quotes in verbatim environments
\IfFileExists{upquote.sty}{\usepackage{upquote}}{}
\IfFileExists{microtype.sty}{% use microtype if available
  \usepackage[]{microtype}
  \UseMicrotypeSet[protrusion]{basicmath} % disable protrusion for tt fonts
}{}
\makeatletter
\@ifundefined{KOMAClassName}{% if non-KOMA class
  \IfFileExists{parskip.sty}{%
    \usepackage{parskip}
  }{% else
    \setlength{\parindent}{0pt}
    \setlength{\parskip}{6pt plus 2pt minus 1pt}}
}{% if KOMA class
  \KOMAoptions{parskip=half}}
\makeatother
\setlength{\emergencystretch}{3em} % prevent overfull lines
\providecommand{\tightlist}{%
  \setlength{\itemsep}{0pt}\setlength{\parskip}{0pt}}
\usepackage{bookmark}
\IfFileExists{xurl.sty}{\usepackage{xurl}}{} % add URL line breaks if available
\urlstyle{same}
\hypersetup{
  hidelinks,
  pdfcreator={LaTeX via pandoc}}

\author{}
\date{}

\begin{document}

\section{Game-Theoretic Framing of the Triadic Balance
Condition}\label{game-theoretic-framing-of-the-triadic-balance-condition}

\textbf{Author:} J. W. Milton, Clarity Coalition

\textbf{Version:} 1.0

\textbf{Date:} 19 May 2026

\textbf{Keywords:} game theory; Nash equilibrium; triadic balance;
tholonic model; classical constants; strategic interaction

\begin{center}\rule{0.5\linewidth}{0.5pt}\end{center}

\subsection{Abstract}\label{abstract}

The tholonic triad assigns three functionally distinct roles
(Negotiation \(N\), Definition \(D\), and Contribution \(C\)) to a
three-variable recurrence whose branches converge to classical
constants: \(\pi/4\), \(\varphi\), \(e\), \(\sqrt{2}\), and \(\ln 2\).
We recast this system as a two-player alternating-move game in which
\(D\) and \(C\) act as strategic agents whose opposing objectives
(constraining versus accumulating) generate a dynamical trajectory for
the payoff state \(N\). The triadic balance condition is identified as
the equilibrium reached when the marginal contributions of \(D\) and
\(C\) to \(N\) are equal and opposite, a condition that reduces to the
fixed-point equation of the underlying recurrence. We prove that this
equilibrium is a pure-strategy Nash equilibrium of the associated stage
game, that the diagonal-invariance and swap-symmetry theorems of the
original framework admit direct characterizations in terms of zero-sum
and symmetric subgame structure, and that the three traversal classes
(Advancing, Self-redefined, Fixed) correspond to distinct information
structures in the game: exogenous shocks, endogenous state evolution,
and constant-strategy play, respectively. The framework provides a new
lens: five classical constants emerge not merely as limits of
recurrences but as equilibrium payoffs of a minimal strategic
interaction.

\begin{center}\rule{0.5\linewidth}{0.5pt}\end{center}

\subsection{1. Introduction}\label{introduction}

Game theory models strategic interaction among agents with distinct
objectives. A central question is how collective outcomes emerge from
individual incentives, particularly when the number of agents is small
and their interests are partially opposed. The triadic framework
introduced by Milton \href{https://tholonia.github.io/the-book/}{Mil24}
and formalized in
\href{https://github.com/baardev/truevalue/blob/main/docnav/Research/papers/1_recursive-tholonic-five-constants/1_recursive-tholonic-five-constants.pdf}{Mil26a}
posits that a three-role partition (a state \(N\) under negotiation, a
bounding force \(D\) that constrains it, and an integrating force \(C\)
that extends it) is the minimal structure sufficient for convergent
recursion toward nontrivial limits. The present paper asks: what does
this system look like when recast as a game?

The translation is natural. The two auxiliary variables \(D\) and \(C\)
pull the state \(N\) in opposite directions: \(D\) contracts, \(C\)
expands. Their interaction can be modeled as a two-player
alternating-move game in which each player selects a magnitude of
influence at each stage, and the state \(N\) evolves as the cumulative
payoff. The balance condition, under which the system converges, becomes
an equilibrium condition in the game.

\textbf{Contributions.} This paper provides:

\begin{enumerate}
\def\labelenumi{\arabic{enumi}.}
\tightlist
\item
  A formal game-theoretic model of the tholonic triad as a two-player
  alternating-move game with an evolving payoff state.
\item
  Identification of the triadic balance condition as a pure-strategy
  Nash equilibrium of the stage game.
\item
  Game-theoretic reinterpretations of the diagonal-invariance and
  swap-symmetry theorems from
  \href{https://github.com/baardev/truevalue/blob/main/docnav/Research/papers/1_recursive-tholonic-five-constants/1_recursive-tholonic-five-constants.pdf}{Mil26a}
  in terms of zero-sum and symmetric subgame structure.
\item
  A mapping of the three traversal classes (Advancing, Self-redefined,
  Fixed) to distinct game-theoretic information structures.
\item
  For each of the five documented branches, an explicit strategic
  interpretation of the players' actions.
\end{enumerate}

\textbf{What this paper does not provide.} New convergence proofs for
the five branches (these are established in
\href{https://github.com/baardev/truevalue/blob/main/docnav/Research/papers/1_recursive-tholonic-five-constants/1_recursive-tholonic-five-constants.pdf}{Mil26a}),
or novel equilibrium existence results beyond those that follow from the
existing recurrence structure. The contribution is the translation
itself: a demonstration that the triadic system admits a coherent and
non-trivial game-theoretic description.

\textbf{Organization.} Section 2 defines the game model. Section 3 gives
strategic interpretations of the three roles. Section 4 analyzes
equilibrium. Section 5 maps traversal classes to game types. Section 6
discusses the five branches in strategic terms. Section 7 covers related
work. Section 8 discusses scope and open questions. Section 9 concludes.

\begin{center}\rule{0.5\linewidth}{0.5pt}\end{center}

\subsection{2. The Triadic Recurrence as a Two-Player
Game}\label{the-triadic-recurrence-as-a-two-player-game}

We consider a two-player alternating-move game \(\mathcal{G}\) played
over stages \(k = 0, 1, 2, \ldots\)

\textbf{Players.} Player Definer (\(\mathcal{D}\)) and Player
Contributor (\(\mathcal{C}\)).

\textbf{State.} A scalar \(N_k \in \mathbb{R}\), the \emph{payoff
state}, initialized at \(N_0\).

\textbf{Actions.} At each stage \(k\), \(\mathcal{D}\) and
\(\mathcal{C}\) each choose a non-negative real parameter: -
\(\mathcal{D}\) selects \(D_k \in \mathbb{R}_{\geq 0}\), the
\emph{defining bound}. - \(\mathcal{C}\) selects
\(C_k \in \mathbb{R}_{\geq 0}\), the \emph{contributing magnitude}.

\textbf{Payoff dynamics.} The state updates according to

\[N_{k+1} = f(N_k; D_k, C_k),\]

where \(f\) is a given update rule. The form of \(f\) defines the game
type. Following
\href{https://github.com/baardev/truevalue/blob/main/docnav/Research/papers/1_recursive-tholonic-five-constants/1_recursive-tholonic-five-constants.pdf}{Mil26a},
we consider branches in which \(f\) is one of:

\[f_{\pi/4}(N; D, C) = N - \frac{1}{D} + \frac{1}{C},\]

\[f_{\varphi}(N; D, C) = D + \frac{D}{N},\]

\[f_{e}(N; D, C) = N + \frac{D}{C},\]

\[f_{\sqrt{2}}(N; D, C) = \frac{N + D/N}{C},\]

\[f_{\ln 2}(N; D, C) = N + (-1)^k \frac{D}{k + C}.\]

\textbf{Strategy spaces.} The strategy of each player is a rule for
selecting \(D_k\) (respectively \(C_k\)) at each stage, possibly
depending on history. Three canonical strategy classes emerge,
corresponding to the traversal classes of
\href{https://github.com/baardev/truevalue/blob/main/docnav/Research/papers/1_recursive-tholonic-five-constants/1_recursive-tholonic-five-constants.pdf}{Mil26a}:

\begin{itemize}
\tightlist
\item
  \textbf{Advancing (Class A):} \(D_{k+1} = D_k + \Delta_D\),
  \(C_{k+1} = C_k + \Delta_C\) for fixed increments
  \(\Delta_D, \Delta_C\). The players receive an exogenous parameter
  injection at each stage.
\item
  \textbf{Self-redefined (Class B):} \(D_{k+1}\) and \(C_{k+1}\) are
  functions of the previous tuple \((N_k, D_k, C_k)\) only. The players
  adapt their actions based on the evolving state.
\item
  \textbf{Fixed (Class C):} \(D_k = D_0\), \(C_k = C_0\) for all \(k\).
  The players commit to constant actions and allow the payoff dynamics
  alone to drive convergence.
\end{itemize}

\textbf{Payoffs.} Each player's per-stage payoff is the marginal effect
of their action on the state:

\[u_{\mathcal{D}}(k) = f(N_k; D_k, C_k) - f(N_k; 0, C_k),\]

\[u_{\mathcal{C}}(k) = f(N_k; D_k, C_k) - f(N_k; D_k, 0),\]

where \(f(N; 0, C)\) denotes the update with \(\mathcal{D}\)'s action
removed (similarly for \(\mathcal{C}\)). The cumulative payoff to each
player is the asymptotic effect on the limit \(N_\infty\).

\textbf{Definition 2.1} (Triadic balance condition). The triad is in
\emph{balance} at stage \(k\) when

\[|u_{\mathcal{D}}(k)| = |u_{\mathcal{C}}(k)|,\]

i.e., the marginal contributions of the two strategic players to the
state change are equal in magnitude. For the five documented branches,
this balance condition is satisfied at the limit point \(N_\infty\).

\begin{center}\rule{0.5\linewidth}{0.5pt}\end{center}

\subsection{3. Strategic Interpretation of the Three
Roles}\label{strategic-interpretation-of-the-three-roles}

The three roles of the triad (Negotiation \(N\), Definition \(D\),
Contribution \(C\)) acquire the following strategic meanings in the
game-theoretic framing.

\subsubsection{\texorpdfstring{3.1 \(N\): The Payoff
State}{3.1 N: The Payoff State}}\label{n-the-payoff-state}

\(N\) is the quantity being \emph{negotiated}. It is not a player with
agency but the object of competition: each player's action modifies
\(N\) in a direction favorable to their objective. \(N\) evolves as the
cumulative result of the tug-of-war between \(\mathcal{D}\) and
\(\mathcal{C}\). In the language of bargaining theory, \(N_k\) is the
provisional agreement after \(k\) rounds of negotiation.

\subsubsection{\texorpdfstring{3.2 \(D\): The Definer
(Constraint)}{3.2 D: The Definer (Constraint)}}\label{d-the-definer-constraint}

\(\mathcal{D}\)'s objective is to \emph{bound} \(N\), to prevent it from
diverging. In the \(\pi/4\) branch, \(\mathcal{D}\) subtracts a term
from \(N\); in the \(\varphi\) branch, \(\mathcal{D}\) supplies the
fixed numerator that defines the contraction; in the \(\sqrt{2}\)
branch, \(\mathcal{D}\) contributes the target value \(2\) toward which
the estimate is pulled; in the \(e\) and \(\ln 2\) branches,
\(\mathcal{D}\) holds the constant numerator \(1\) that bounds each
term's contribution.

\(\mathcal{D}\) can be understood as a \emph{minimizing} player in a
zero-sum component: its action reduces (or bounds) the state relative to
what \(\mathcal{C}\) alone would produce.

\subsubsection{\texorpdfstring{3.3 \(C\): The Contributor
(Accumulation)}{3.3 C: The Contributor (Accumulation)}}\label{c-the-contributor-accumulation}

\(\mathcal{C}\)'s objective is to \emph{extend} and \emph{integrate}. In
the \(\pi/4\) branch, \(\mathcal{C}\) adds a positive term to \(N\),
countering \(\mathcal{D}\)'s subtraction. In the \(e\) branch,
\(\mathcal{C}\) expands the denominator factorially, controlling the
convergence rate. In the \(\sqrt{2}\) branch, \(\mathcal{C}\) provides
the averaging divisor that synthesizes the correction. In the \(\ln 2\)
branch, \(\mathcal{C}\) provides the harmonic offset in the denominator,
structuring the alternating series.

\(\mathcal{C}\) can be understood as a \emph{maximizing} or
\emph{accumulating} player, extending the state in a direction opposite
to \(\mathcal{D}\).

\subsubsection{3.4 Irreducibility as Strategic
Necessity}\label{irreducibility-as-strategic-necessity}

Lemma 3.1 of
\href{https://github.com/baardev/truevalue/blob/main/docnav/Research/papers/1_recursive-tholonic-five-constants/1_recursive-tholonic-five-constants.pdf}{Mil26a}
(triadic irreducibility) acquires a game-theoretic interpretation: a
single-player game with only a state variable produces either unbounded
growth or trivial convergence. Two players with opposing incentives are
the minimum structure for non-trivial equilibrium; the state \(N\) is
the arena in which their contest plays out. Fewer than three roles give
either no competition (one player) or competition without a scoreboard
(two players, no state). The triad as a two-player game with payoff
state is thus the minimal strategic structure for convergent dynamics to
a non-trivial limit.

\begin{center}\rule{0.5\linewidth}{0.5pt}\end{center}

\subsection{4. Equilibrium Analysis}\label{equilibrium-analysis}

\subsubsection{4.1 The Balance Condition as Nash
Equilibrium}\label{the-balance-condition-as-nash-equilibrium}

Consider the stage game at iteration \(k\), holding the history
\((N_0, \ldots, N_{k-1})\) fixed. Players \(\mathcal{D}\) and
\(\mathcal{C}\) choose \(D_k\) and \(C_k\) simultaneously (or,
equivalently, in alternating order with stage-complete information). The
per-stage payoff change is

\[\Delta N_k = f(N_k; D_k, C_k) - N_k.\]

\textbf{Definition 4.1} (Stage-game Nash equilibrium). A pair
\((D_k^*, C_k^*)\) is a pure-strategy Nash equilibrium of the stage game
if, for all admissible \(D_k, C_k\),

\[f(N_k; D_k^*, C_k^*) - f(N_k; D_k, C_k^*) \leq 0,\]

\[f(N_k; D_k^*, C_k^*) - f(N_k; D_k^*, C_k) \leq 0,\]

where the inequalities reflect the minimizing objective of
\(\mathcal{D}\) and the maximizing objective of \(\mathcal{C}\),
respectively (sign conventions differ by branch; the general condition
is that neither player can unilaterally improve its per-stage marginal
contribution).

\textbf{Proposition 4.2} (Balance implies equilibrium for Class A). For
the Class A update \(f(N; D, C) = N - 1/D + 1/C\) with fixed
\(D, C > 0\), the unique stage-game Nash equilibrium is any \((D, C)\)
with \(D = C\). At this equilibrium, \(\Delta N = 0\) and the balance
condition \(|u_{\mathcal{D}}| = |u_{\mathcal{C}}|\) holds trivially.

\emph{Proof.} \(\mathcal{D}\) chooses \(D\) to minimize
\(N - 1/D + 1/C\); equivalently, to maximize \(1/D\). Since \(1/D\) is
decreasing in \(D\), \(\mathcal{D}\) prefers larger \(D\).
\(\mathcal{C}\) chooses \(C\) to maximize \(N - 1/D + 1/C\);
equivalently, to maximize \(1/C\), which is achieved for smaller \(C\).
With \(D = C\), the terms cancel: \(\Delta N = 0\). Neither player can
achieve a net positive contribution unilaterally without the other's
response. This is the unique symmetric equilibrium of the one-shot stage
game. \(\square\)

The diagonal-invariance theorem of
\href{https://github.com/baardev/truevalue/blob/main/docnav/Research/papers/1_recursive-tholonic-five-constants/1_recursive-tholonic-five-constants.pdf}{Mil26a}
(Proposition 6.2) states that \(D = C\) yields no net state change;
Proposition 4.2 shows this is not merely an algebraic identity but an
equilibrium property.

\textbf{Proposition 4.3} (Swap symmetry as zero-sum subgame). For the
Class A update, let \(\mathcal{G}_{(D_0, C_0)}\) denote the game
initialized with \((D_0, C_0)\). Then

\[N_\infty^{\mathcal{G}_{(D_0, C_0)}} + N_\infty^{\mathcal{G}_{(C_0, D_0)}} = 2 N_0.\]

The two games are strategic mirrors: swapping the players' initial
actions reverses the sign of the net contribution at every stage, making
the combined payoff constant. This is precisely the zero-sum property
for the \emph{difference} between the two game instances.

\emph{Proof.} Follows from Proposition 6.3 of
\href{https://github.com/baardev/truevalue/blob/main/docnav/Research/papers/1_recursive-tholonic-five-constants/1_recursive-tholonic-five-constants.pdf}{Mil26a}.
The strategic interpretation is that \(\mathcal{G}_{(D_0, C_0)}\) and
\(\mathcal{G}_{(C_0, D_0)}\) form a zero-sum pair: the gain of
\(\mathcal{C}\) in one game is exactly the loss of \(\mathcal{D}\) in
the mirrored game, and vice versa. \(\square\)

\subsubsection{4.2 The Limit as Correlated
Equilibrium}\label{the-limit-as-correlated-equilibrium}

In the full dynamic game, players do not choose \(D_k\) and \(C_k\)
myopically; they follow a prescribed strategy profile (the traversal
rule of the branch). The trajectory \((N_k, D_k, C_k)_{k \geq 0}\)
converges to a limit \(N_\infty\) at which the marginal contributions of
\(\mathcal{D}\) and \(\mathcal{C}\) vanish asymptotically:

\[\lim_{k \to \infty} \bigl| f(N_k; D_k, C_k) - N_k \bigr| = 0.\]

This limit point can be interpreted as a \emph{correlated equilibrium}
of the repeated game: the prescribed strategy profile (traversal rule)
is a correlation device under which neither player has a profitable
unilateral deviation that would alter the asymptotic payoff, given that
the other player adheres to the profile. The five classical constants
emerge as the equilibrium payoffs of five distinct correlated
equilibria.

\begin{center}\rule{0.5\linewidth}{0.5pt}\end{center}

\subsection{5. Traversal Classes as Information
Structures}\label{traversal-classes-as-information-structures}

The three traversal classes of
\href{https://github.com/baardev/truevalue/blob/main/docnav/Research/papers/1_recursive-tholonic-five-constants/1_recursive-tholonic-five-constants.pdf}{Mil26a}
map to three distinct game-theoretic information structures.

\subsubsection{5.1 Class A (Advancing): Exogenous
Shocks}\label{class-a-advancing-exogenous-shocks}

In the advancing class, \(D_k\) and \(C_k\) each receive a fixed
external increment at every stage. In game-theoretic terms, this is a
game with \emph{exogenous state variables}: a component of each player's
action is determined by an outside process (the geometric axis
structure, parameterized by \(\Delta = 4\)). The players' strategic
choice is only the \emph{initial seed}; thereafter, the environment
dictates the increment. This structure is reminiscent of stochastic
games with a deterministic external driver {[}Sha53{]}, or of repeated
games with a publicly observed state that evolves independently of
actions.

The \(\pi/4\) branch is the sole Class A representative. The step size
\(\Delta = 4\) is derived from the Instantiation axis multiplier
\(\mathrm{istep} = 2\) via \(d_{\mathrm{step}} = \mathrm{istep}^2 = 4\)
and \(c_{\mathrm{step}} = 2 \times \mathrm{istep} = 4\), yielding a
single parameter-free increment. The initial seeds
\((N_0, D_0, C_0) = (1, 3, 5)\) are themselves taken from the axis
multipliers of the underlying geometric configuration.

\subsubsection{5.2 Class B (Self-redefined): Endogenous State
Evolution}\label{class-b-self-redefined-endogenous-state-evolution}

In the self-redefined class, \(D_k\) and \(C_k\) evolve as functions of
the previous tuple only. In game-theoretic terms, this is a game of
\emph{complete information with endogenous state}: each player
conditions its next-period action on the full history of play, and the
state transitions are deterministic functions of that history. This
structure is closely related to \emph{dynamic games with Markov perfect
equilibria} {[}MS94{]}, where strategies depend only on the
payoff-relevant state.

The \(\varphi\) and \(e\) branches belong to Class B. In the \(\varphi\)
branch, the Fibonacci swap \((D, C) \leftarrow (C, C + D)\) is a
deterministic state transition rule that does not depend on any external
signal. In the \(e\) branch, \(C_k = k!\) is a deterministic function of
the stage index, which is itself part of the endogenous game state.

\subsubsection{5.3 Class C (Fixed): Constant-Strategy
Play}\label{class-c-fixed-constant-strategy-play}

In the fixed class, \(D_k = D_0\) and \(C_k = C_0\) for all \(k\). In
game-theoretic terms, this is a game in which both players commit to
\emph{stationary strategies}: their actions are time-invariant.
Convergence is driven entirely by the form of the payoff function \(f\)
and the accumulation over stages. This is the simplest strategic
structure: a repeated game with constant actions, where the only
strategic choice is the initial commitment.

The \(\sqrt{2}\) and \(\ln 2\) branches belong to Class C. Despite the
identical structure at the strategy level, the payoff functions differ:
the \(\sqrt{2}\) branch uses a Newton-type contraction, while the
\(\ln 2\) branch uses an alternating series. Both produce non-trivial
limits from constant actions, demonstrating that stationary strategies
need not yield degenerate outcomes when the payoff function itself
encodes sufficient structure.

\begin{center}\rule{0.5\linewidth}{0.5pt}\end{center}

\subsection{6. Strategic Interpretation of the Five
Branches}\label{strategic-interpretation-of-the-five-branches}

Each of the five documented branches of the tholonic ladder
\href{https://github.com/baardev/truevalue/blob/main/docnav/Research/papers/1_recursive-tholonic-five-constants/1_recursive-tholonic-five-constants.pdf}{Mil26a}
admits a distinct strategic narrative under the game-theoretic framing.

\subsubsection{\texorpdfstring{6.1 \(\pi/4\): The Alternating
Tug-of-War}{6.1 \textbackslash pi/4: The Alternating Tug-of-War}}\label{pi4-the-alternating-tug-of-war}

\(\mathcal{D}\) and \(\mathcal{C}\) alternately subtract and add terms
of the form \(1/(4n-1)\) and \(1/(4n+1)\), with both players' actions
advancing in lockstep by \(\Delta = 4\). This is an
\emph{alternating-move bargaining game} with an exogenous clock. The
balance condition is dynamic: at each stage the two contributions
partially cancel, and only in the limit do they settle to \(\pi/4\). The
equilibrium is correlated: neither player can accelerate or decelerate
the clock, so deviation is impossible without altering the rules of the
game.

\subsubsection{\texorpdfstring{6.2 \(\varphi\): Golden
Bargaining}{6.2 \textbackslash varphi: Golden Bargaining}}\label{varphi-golden-bargaining}

\(\mathcal{D}\) supplies the fixed unit \(1\) as the definitional
baseline; \(\mathcal{C}\) reciprocates with the evolving ratio \(D/N\),
a Fibonacci-weighted accumulation. The players swap roles at each stage
(the swap symmetry \(D \leftarrow C\)), creating a \emph{reciprocal
bargaining protocol} in which each player's action becomes the other's
baseline in the next round. The fixed point \(\varphi = 1 + 1/\varphi\)
is the equilibrium of this reciprocal dynamic: the point at which the
exchange stabilizes.

\subsubsection{\texorpdfstring{6.3 \(e\): Factorial
Escalation}{6.3 e: Factorial Escalation}}\label{e-factorial-escalation}

\(\mathcal{D}\) commits to a constant action (\(D_k = 1\)) throughout.
\(\mathcal{C}\) escalates its denominator factorially (\(C_k = k!\)),
rapidly diluting the per-stage contribution. The equilibrium payoff
\(e\) is the limit of this one-sided escalation. In strategic terms,
this is a \emph{war of attrition} in reverse: rather than escalating
costs, \(\mathcal{C}\) escalates the denominator, making each successive
contribution vanishingly small. \(\mathcal{D}\)'s passivity (constant
\(D = 1\)) is the stable response to \(\mathcal{C}\)'s factorial
strategy.

\subsubsection{\texorpdfstring{6.4 \(\sqrt{2}\): The Babylonian
Standoff}{6.4 \textbackslash sqrt\{2\}: The Babylonian Standoff}}\label{sqrt2-the-babylonian-standoff}

Both players commit to constant actions (\(D_0 = C_0 = 2\)). The payoff
function \(f(N) = (N + 2/N)/2\) is a contraction mapping with
\(\sqrt{2}\) as its unique fixed point. The game is a
\emph{fixed-commitment game} in which the strategic content is entirely
in the initial choice of target (\(D = 2\)) and averaging weight
(\(C = 2\)). Once committed, the players have no further moves;
convergence is guaranteed by the contraction property of the payoff
function.

\subsubsection{\texorpdfstring{6.5 \(\ln 2\): The Harmonic
Alternation}{6.5 \textbackslash ln 2: The Harmonic Alternation}}\label{ln-2-the-harmonic-alternation}

Both players commit to constant actions (\(D_0 = C_0 = 1\)). The payoff
function injects an alternating sign \((-1)^k\), creating a
\emph{turn-based zero-sum structure within a single payoff function}.
The players themselves do not alternate; the alternation is encoded in
the game rules. The limit \(\ln 2\) is the asymptotic value of the
alternating harmonic series, interpretable as the \emph{fair division}
outcome of an infinite sequence of alternating positive and negative
contributions.

\begin{center}\rule{0.5\linewidth}{0.5pt}\end{center}

\subsection{7. Related Work}\label{related-work}

\textbf{Evolutionary game theory and replicator dynamics.} The
convergence of the triad to a limit point resembles the convergence of
population proportions under replicator dynamics to an evolutionarily
stable strategy {[}HS98, Wei95{]}. In both settings, a dynamical system
defined by pairwise interactions tends toward a fixed point. The
difference is that replicator dynamics operate over a simplex of
population proportions, while the triad operates over a three-role
partition with asymmetric functional roles. Whether each branch can be
embedded as a projection of a higher-dimensional replicator system is an
open question.

\textbf{Bargaining theory.} The sequential nature of the triad
(alternating contributions from \(D\) and \(C\)) is structurally similar
to the Rubinstein alternating-offers bargaining model {[}Rub82{]}. In
the Rubinstein model, two players alternate in proposing divisions of a
surplus; the unique subgame-perfect equilibrium converges to a split
determined by discount factors. In the triad, \(D\) and \(C\) alternate
in modifying \(N\), and the limit is determined by the traversal rule. A
formal reduction of any branch to a Rubinstein-type game with
appropriately chosen discount factors would strengthen the connection.

\textbf{Potential games.} Monderer and Shapley {[}MS96{]} introduced
potential games, in which the incentive of each player to change
strategy can be expressed by a single global potential function. If the
triad update \(f(N; D, C)\) can be written as the gradient of a
potential \(\Phi(N, D, C)\), then the convergence to equilibrium follows
from the general theory of potential games. For the \(\sqrt{2}\) branch,
the potential is \(\Phi(N) = (N - \sqrt{2})^2\), whose gradient descent
recovers the Newton iteration. Characterizing the potential for the
remaining branches is a direction for future work.

\textbf{Learning in games.} Fudenberg and Levine {[}FL98{]} study
learning dynamics in which players adjust strategies based on past
payoffs. The triad can be viewed as a learning system in which \(D\) and
\(C\) adjust their influence parameters over time, with \(N\) encoding
the cumulative result of this learning. The three traversal classes then
correspond to three learning paradigms: exogenous curriculum (Class A),
self-supervised adaptation (Class B), and fixed-policy evaluation (Class
C).

\textbf{Zero-sum and symmetric games.} The swap-symmetry theorem
(Proposition 4.3 above) identifies a zero-sum structure in the Class A
framework. Von Neumann and Morgenstern {[}vNM44{]} established the
theory of zero-sum two-player games; the triad extends this structure to
a three-component system in which the zero-sum property holds between
\emph{paired game instances} rather than within a single instance.

\begin{center}\rule{0.5\linewidth}{0.5pt}\end{center}

\subsection{8. Discussion}\label{discussion}

\textbf{Is the game-theoretic framing merely relabeling?} A fair
question. The algebraic content of the three-variable recurrences is
unchanged. What the game-theoretic translation adds is:

\begin{enumerate}
\def\labelenumi{\arabic{enumi}.}
\item
  \textbf{Causal interpretation.} \(D\) and \(C\) are not merely
  parameters but agents with opposing incentives. The fact that their
  marginal contributions cancel at the limit is not a coincidence; it
  reflects the equilibrium of a minimal strategic interaction. The
  balance condition is the point at which neither agent's unilateral
  deviation would advance its objective.
\item
  \textbf{Taxonomic unification.} The three traversal classes map
  cleanly to three canonical game-theoretic information structures
  (exogenous, endogenous, stationary). This mapping is not forced by the
  recurrence formulation; it is a discovered correspondence that adds a
  layer of explanation for why the three classes behave differently.
\item
  \textbf{Connective tissue.} The translation bridges the tholonic
  framework to existing game-theoretic literature (bargaining, potential
  games, learning in games, zero-sum theory). These connections provide
  vocabulary and theorems that may prove useful for the open conjectures
  of
  \href{https://github.com/baardev/truevalue/blob/main/docnav/Research/papers/1_recursive-tholonic-five-constants/1_recursive-tholonic-five-constants.pdf}{Mil26a},
  particularly Conjecture 8.3 (finite classification of named limits).
\end{enumerate}

\textbf{Open questions.} Several questions raised by the game-theoretic
framing remain open:

\begin{enumerate}
\def\labelenumi{\arabic{enumi}.}
\tightlist
\item
  Can each of the five branches be formally reduced to a Rubinstein
  bargaining game with a specific discount factor?
\item
  Does a single potential function \(\Phi(N, D, C)\) exist whose
  gradient descent recovers all five branches?
\item
  Does the game-theoretic perspective suggest new branches (i.e., new
  equilibrium limits) that were not visible from the recurrence
  perspective alone?
\item
  Can the convergence-rate hierarchy (\(O(1/k)\) for \(\pi/4\) and
  \(\ln 2\), faster for the rest) be derived from the information
  structure of the corresponding game class (Class A/Class C slow, Class
  B fast) as a general theorem?
\end{enumerate}

\textbf{Limitations.} This paper provides a descriptive translation, not
a predictive extension. The game-theoretic model does not (yet) generate
new constants or new convergence guarantees beyond those already
established by the recurrence framework. The value is in the
reinterpretation: a lens through which the structural theorems of
\href{https://github.com/baardev/truevalue/blob/main/docnav/Research/papers/1_recursive-tholonic-five-constants/1_recursive-tholonic-five-constants.pdf}{Mil26a}
acquire strategic meaning.

\begin{center}\rule{0.5\linewidth}{0.5pt}\end{center}

\subsection{9. Conclusion}\label{conclusion}

We have recast the tholonic triad \((N, D, C)\) as a two-player
alternating-move game in which the Definer (\(\mathcal{D}\)) and the
Contributor (\(\mathcal{C}\)) act as strategic agents with opposing
objectives, and the state \(N\) evolves as their cumulative payoff. The
triadic balance condition, defined as equality of the marginal
contributions of the two players, is shown to be a pure-strategy Nash
equilibrium of the stage game for the Class A update, and a correlated
equilibrium of the repeated game across all five branches.

The swap-symmetry and diagonal-invariance theorems of the original
recurrence framework acquire direct game-theoretic characterizations as
zero-sum and symmetric equilibrium properties. The three traversal
classes map to three canonical game-theoretic information structures:
exogenous shocks (Class A), endogenous state evolution with complete
information (Class B), and stationary strategies (Class C). Each of the
five converged constants (\(\pi/4\), \(\varphi\), \(e\), \(\sqrt{2}\),
\(\ln 2\)) is reinterpreted as the equilibrium payoff of a distinct
strategic interaction.

This translation does not alter the mathematical content of the
recurrence framework, but it supplies a new vocabulary. The minimal
strategic structure that underlies the triad (two opposing agents plus a
negotiated state) connects naturally to bargaining theory, potential
games, and learning dynamics. Whether this connection can be leveraged
to resolve the open conjectures of
\href{https://github.com/baardev/truevalue/blob/main/docnav/Research/papers/1_recursive-tholonic-five-constants/1_recursive-tholonic-five-constants.pdf}{Mil26a}
remains to be seen.

\begin{center}\rule{0.5\linewidth}{0.5pt}\end{center}

\subsection{References}\label{references}

{[}FL98{]} D. Fudenberg and D. K. Levine, \emph{The Theory of Learning
in Games}, MIT Press, 1998.

{[}HS98{]} J. Hofbauer and K. Sigmund, \emph{Evolutionary Games and
Population Dynamics}, Cambridge University Press, 1998.

\href{https://tholonia.github.io/the-book/}{Mil24} J. W. Milton,
\emph{Tholonia: The Existential Mechanics of Awareness}, self-published,
2024. Available at https://tholonia.github.io/the-book/.

\href{https://github.com/baardev/truevalue/blob/main/docnav/Research/papers/1_recursive-tholonic-five-constants/1_recursive-tholonic-five-constants.pdf}{Mil26a}
J. W. Milton, ``Emergence of Classical Constants from a Minimal
Recursive Triadic Framework,'' arXiv preprint, 2026.

{[}MS94{]} E. Maskin and J. Tirole, ``Markov Perfect Equilibrium,''
\emph{Journal of Economic Theory}, 1994.

{[}MS96{]} D. Monderer and L. S. Shapley, ``Potential Games,''
\emph{Games and Economic Behavior}, 14(1), 1996, pp.~124--143.

{[}Rub82{]} A. Rubinstein, ``Perfect Equilibrium in a Bargaining
Model,'' \emph{Econometrica}, 50(1), 1982, pp.~97--109.

{[}Sha53{]} L. S. Shapley, ``Stochastic Games,'' \emph{Proceedings of
the National Academy of Sciences}, 39(10), 1953, pp.~1095--1100.

{[}vNM44{]} J. von Neumann and O. Morgenstern, \emph{Theory of Games and
Economic Behavior}, Princeton University Press, 1944.

{[}Wei95{]} J. W. Weibull, \emph{Evolutionary Game Theory}, MIT Press,
1995.

\begin{center}\rule{0.5\linewidth}{0.5pt}\end{center}

\subsection{Appendix: Figure assets}\label{appendix-figure-assets}

All figures for this manuscript use filenames under
\texttt{papers/figures/} with the prefix \textbf{\texttt{4\_}} (for
example \texttt{figures/4\_stage-game-payoffs.png}). This version of the
preprint contains no embedded figures; any added in revision must follow
the same prefix so paths stay unique across papers in this repository.

\end{document}
